The Thayer-Direct result shows that the controlled synthetic task is learnable:
learned reconstruction substantially improves affected-region error over
identity and threshold baselines. The hard stress test shows that normal
held-out performance is not enough to characterize robustness.

Thayer-Residual helps because the model learns contaminant light to subtract
rather than redrawing the whole target. This can preserve target structure in
regions that are already correct in the blended image. Thayer-BR v0.1 helps
because core overlap and bright/similar-size contaminants are
sampled deliberately during optimization instead of appearing only by chance.

The result should not be overstated. Thayer-Direct and Thayer-Residual still
win on some individual samples. Some failures involve ambiguous overlap,
target-core obstruction, over-smoothing, or lost target detail. Also, high blend
severity is not identical to model difficulty: an obvious contaminant can be
severe but easy to subtract, while a small overlap can be hard if it destroys
central target structure.

The current benchmark remains synthetic. More realistic sky backgrounds, PSF
variation, correlated source environments, and survey-like noise are needed
before making broader astronomical claims.
